\begin{table}[!htbp]\centering
\caption{The suicide null is not an artifact of network-spillover-contaminated controls}
\label{tab:network-spillover}
\small
\begin{tabular}{lrrrr}
\toprule
Control set & Suicide ATT & 95\% CI & Controls dropped & $N$ controls \\
\midrule
All controls (headline) & +0.28 & [-0.72, +1.29] & 0 & 5{,}461 \\
\midrule
Drop co-patiency weight $\geq 0.05$ & +0.35 & [-0.69, +1.39] & 1{,}173 & 4{,}288 \\
Drop co-patiency weight $\geq 0.10$ & +0.36 & [-0.67, +1.39] & 1{,}035 & 4{,}426 \\
Drop co-patiency weight $\geq 0.25$ & +0.34 & [-0.68, +1.36] & 777 & 4{,}684 \\
Drop co-patiency weight $\geq 0.50$ & +0.33 & [-0.69, +1.35] & 589 & 4{,}872 \\
Drop co-patiency weight $\geq 1.00$ & +0.32 & [-0.69, +1.33] & 364 & 5{,}097 \\
\bottomrule
\multicolumn{5}{p{0.94\textwidth}}{\footnotesize Notes: Population-weighted staggered Sun-Abraham event study for suicide mortality per $100{,}000$, municipality and year fixed effects, standard errors clustered by municipality. A never-treated municipality is treated as network-contaminated if it shares a co-patiency edge of weight at least $\theta$ with any exposed catchment in the projected municipality network (municipalities that co-use the same hospitals are exposed to the same closure-induced rerouting). Because spillover attenuates a true effect toward zero, dropping contaminated controls would \emph{raise} the estimate if an effect were hidden; instead the ATT stays a bounded null across thresholds. This complements the shared-hub, shared-substitute, and flow-similarity contamination checks in the main robustness battery.}\\
\end{tabular}
\end{table}

